\documentclass[11pt]{article}
\usepackage[a4paper,margin=2.5cm]{geometry}
\usepackage{amsmath,amssymb,amsthm}
\usepackage[british]{babel}
\newtheorem{proposition}{Proposition}
\newtheorem{lemma}[proposition]{Lemma}
\theoremstyle{remark}
\newtheorem{remark}[proposition]{Remark}
\newcommand{\E}{\mathbb{E}}
\newcommand{\Q}{\mathbb{Q}}
\renewcommand{\P}{\mathbb{P}}

\title{Crash insurance and the G10 carry premium: theory notes}
\author{Amanjeet Singh}
\date{Version of 24 September 2026}

\begin{document}
\maketitle

\noindent These notes prove the results listed in \texttt{theory/README.md}. Results that are standard are marked as such and attributed; the proofs written here are the author's own unless a source is named. Throughout, $\Phi$ and $\varphi$ are the standard normal distribution and density functions, $s=\sigma\sqrt{\tau}$, and for forward $F$ and strike $K$
\[
d_\pm(K)=\frac{\ln(F/K)\pm s^2/2}{s}.
\]
The identity $K\varphi(d_-)=F\varphi(d_+)$ is used repeatedly.

\section*{R1. Domestic and foreign risk-neutral measures}

This is the change of numeraire of Geman, El Karoui and Rochet (1995); the statement and proof are included for completeness. Girsanov's theorem and It\^o's formula are used in the form given by Shreve (2004).

Let $(\Omega,\mathcal F,(\mathcal F_t)_{t\in[0,T]},\Q^d)$ satisfy the usual conditions, let $B^d,B^f$ be strictly positive adapted money-market accounts with $B^d_0=B^f_0=1$, and let $S>0$ be adapted and c\`adl\`ag. Assume $Z_t=S_tB^f_t/(S_0B^d_t)$ is a true $\Q^d$-martingale and define $d\Q^f/d\Q^d=Z_T$ on $\mathcal F_T$.

\begin{proposition}
(a) $\Q^f$ is a probability measure equivalent to $\Q^d$. (b) For $\mathcal F_T$-measurable $X\ge0$ or $X\in L^1(\Q^f)$, $\E^{\Q^f}[X\mid\mathcal F_t]=\E^{\Q^d}[Z_TX\mid\mathcal F_t]/Z_t$. (c) An adapted process $Y$ with $Y_t\in L^1(\Q^f)$ is a $\Q^f$-martingale if and only if $ZY$ is a $\Q^d$-martingale; in particular a foreign price $V^f$ satisfies: $V^f/B^f$ is a $\Q^f$-martingale if and only if $SV^f/B^d=S_0Z\,(V^f/B^f)$ is a $\Q^d$-martingale. (d) If $dS_t/S_t=(r_d-r_f)\,dt+\sigma\,dW_t$ under $\Q^d$ with constant coefficients, then $W^f_t=W_t-\sigma t$ is a $\Q^f$-Brownian motion, $d(1/S_t)=(1/S_t)[(r_f-r_d)\,dt-\sigma\,dW^f_t]$, and $\E^{\Q^d}S_T=F$, $\E^{\Q^f}(1/S_T)=1/F$ with $F=S_0e^{(r_d-r_f)T}$.
\end{proposition}

\begin{proof}
(a) $Z_T>0$ and $\E^{\Q^d}Z_T=Z_0=1$, so $\Q^f$ is a probability measure; positivity of $Z_T$ gives equivalence. (b) For $\mathcal F_t$-measurable $Y\ge0$, the tower property and the martingale property give $\E^{\Q^f}Y=\E^{\Q^d}[Z_TY]=\E^{\Q^d}[Z_tY]$. For $A\in\mathcal F_t$,
\[
\E^{\Q^f}[X\mathbf 1_A]=\E^{\Q^d}\bigl[\E^{\Q^d}[Z_TX\mid\mathcal F_t]\mathbf 1_A\bigr]=\E^{\Q^f}\Bigl[\frac{\E^{\Q^d}[Z_TX\mid\mathcal F_t]}{Z_t}\mathbf 1_A\Bigr],
\]
which characterises the conditional expectation. (c) Apply (b) on $[s,t]$: $\E^{\Q^f}[Y_t\mid\mathcal F_s]=\E^{\Q^d}[Z_tY_t\mid\mathcal F_s]/Z_s$, so $Y$ is a $\Q^f$-martingale exactly when $\E^{\Q^d}[Z_tY_t\mid\mathcal F_s]=Z_sY_s$. (d) Here $Z_t=\exp(\sigma W_t-\sigma^2t/2)$, so Girsanov's theorem gives the first claim, and It\^o's formula for $1/S$ gives the dynamics. The expectations follow from the lognormal law under each measure.
\end{proof}

\begin{remark}
Jensen's inequality gives $\E^{\Q^d}(1/S_T)>1/\E^{\Q^d}S_T=1/F$. There is no contradiction with (d): the foreign investor's forward price is an expectation under $\Q^f$, not $\Q^d$ (Siegel's paradox). If $Z$ is only a strict local martingale, then $\E Z_T<1$ and $Z_T$ does not define an equivalent probability measure. Carr, Fisher and Ruf (2014) construct the foreign measure in that case as a F\"ollmer measure that is not equivalent to $\Q^d$, and restore the domestic--foreign symmetry with a modified pricing operator.
\end{remark}

\section*{R2. Arbitrage-free call prices}

The link between the second strike derivative of call prices and state prices is due to Breeden and Litzenberger (1978). The characterisation below, with its measure-theoretic proof, is stated and proved here.

\begin{proposition}
Let $c:[0,\infty)\to\mathbb{R}$ and $F>0$. There is a random variable $X\ge 0$ with $\E X=F$ and $c(K)=\E(X-K)^+$ for all $K\ge0$ if and only if (i) $c$ is convex, (ii) $c(0)=F$, (iii) $c'_+(0)\ge -1$ and (iv) $c(K)\to0$ as $K\to\infty$. The law of $X$ is then unique, with distribution function $G=1+c'_+$ on $[0,\infty)$ and $\P(X=0)=1+c'_+(0)$.
\end{proposition}

\begin{proof}
\emph{Necessity.} Convexity of $c$ follows from convexity of $K\mapsto(x-K)^+$ and linearity of expectation, and $c(0)=\E X=F$. Since $0\le(X-K)^+\le X\in L^1$ and $(X-K)^+\to0$ pointwise, dominated convergence gives (iv). For $h>0$, the difference quotient $[(X-K-h)^+-(X-K)^+]/h$ is bounded by $1$ in absolute value and converges to $-\mathbf 1\{X>K\}$ as $h\downarrow0$, so $c'_+(K)=-\P(X>K)$. In particular $c'_+(0)=-\P(X>0)\ge-1$.

\emph{Sufficiency.} A convex function on $[0,\infty)$ has a right derivative $c'_+$ that is nondecreasing and right-continuous, and $c$ is absolutely continuous on compact intervals with $c(b)-c(a)=\int_a^b c'_+(u)\,du$. If $c'_+(K_0)>0$ for some $K_0$, then $c(K)\ge c(K_0)+c'_+(K_0)(K-K_0)\to\infty$, contradicting (iv); so $c'_+\le0$, $c$ is nonincreasing and, by (iv), $c\ge0$. Let $L=\lim_{K\to\infty}c'_+(K)\le0$. Since $c'_+$ is nondecreasing, $c'_+\le L$ everywhere, so $c(K)\le F+LK$; $L<0$ would force $c<0$ for large $K$. Hence $L=0$.

Define $G(K)=0$ for $K<0$ and $G(K)=1+c'_+(K)$ for $K\ge0$. Then $G$ is nondecreasing and right-continuous, $G(0)=1+c'_+(0)\ge0$ by (iii), and $G(K)\to1$. It is therefore the distribution function of a random variable $X\ge0$ (Lebesgue--Stieltjes construction). By Tonelli's theorem,
\[
\E(X-K)^+=\int_K^\infty \P(X>u)\,du=\int_K^\infty\bigl(1-G(u)\bigr)du=-\int_K^\infty c'_+(u)\,du=c(K)-\lim_{b\to\infty}c(b)=c(K),
\]
and $\E X=c(0)=F$. Uniqueness holds because $c$ determines $c'_+$ and hence $G$.
\end{proof}

\begin{remark}
Extending $c$ by $c(K)=F-K$ for $K<0$, the second distributional derivative of $c$ on $\mathbb{R}$ is the law of $X$; the jump of $c'$ at $0$ carries the atom $\P(X=0)$. On a finite strike grid the necessary conditions reduce to slopes in $[-1,0]$ that are nondecreasing, which is the discrete test used in \texttt{src/qef/fx/arbitrage.py}. The implementation also applies the density condition on implied total variance of Gatheral and Jacquier (2014, eq.~(2.1) and Lemma~2.2).
\end{remark}

\section*{R5(b). Premium-adjusted call delta}

The premium-adjusted delta, its non-monotonicity in the strike and the market practice of taking the strike to the right of the delta maximum are described by Reiswich and Wystup (2012). The proposition below makes the shape precise and is proved here.

Let $D\in(0,1]$ and $h(K)=D\,(K/F)\,\Phi\bigl(d_-(K)\bigr)$ for $K>0$. With spot delta, $D$ is the base-currency discount factor; with forward delta, $D=1$.

\begin{lemma}
The inverse Mills ratio $\lambda(d)=\varphi(d)/\Phi(d)$ is strictly decreasing on $\mathbb{R}$, with $\lambda(d)\to\infty$ as $d\to-\infty$ and $\lambda(d)\to0$ as $d\to\infty$.
\end{lemma}

\begin{proof}
Differentiation gives $\lambda'(d)=-\lambda(d)\,[d+\lambda(d)]$, so it suffices that $d\,\Phi(d)+\varphi(d)>0$. This function has derivative $\Phi(d)>0$ and tends to $0$ as $d\to-\infty$; hence $d\,\Phi(d)+\varphi(d)=\int_{-\infty}^d\Phi(u)\,du>0$. The limits follow from $\Phi\to0$, $\varphi\to0$ and the Mills-ratio asymptotics $\Phi(d)\sim\varphi(d)/|d|$ as $d\to-\infty$. Inequalities of this type for the Mills ratio are classical; see Sampford (1953).
\end{proof}

\begin{proposition}
The function $h$ is continuous on $(0,\infty)$, tends to $0$ as $K\downarrow0$ and as $K\to\infty$, and satisfies $h<D$. There is a unique $K^*$ such that $h$ is strictly increasing on $(0,K^*]$ and strictly decreasing on $[K^*,\infty)$; $K^*=F\exp(-d^*s-s^2/2)$, where $d^*$ is the unique solution of $\lambda(d)=s$. Each $\Delta\in(0,h(K^*))$ has exactly two preimages, one on each side of $K^*$.
\end{proposition}

\begin{proof}
As $K\downarrow0$, $d_-\to\infty$, so $h\to0$. As $K\to\infty$, $d_-\to-\infty$ and, by the Mills-ratio asymptotics and the identity above, $K\Phi(d_-)\sim K\varphi(d_-)/|d_-|=F\varphi(d_+)/|d_-|\to0$. Moreover $F\Phi(d_+)-K\Phi(d_-)$ is the undiscounted call price, which is positive, so $h<D\,\Phi(d_+)<D$.

Since $d_-'(K)=-1/(Ks)$,
\[
h'(K)=\frac{D}{F}\Bigl[\Phi(d_-)-\frac{\varphi(d_-)}{s}\Bigr]=\frac{D\,\Phi(d_-)}{F\,s}\,\bigl[s-\lambda(d_-)\bigr].
\]
By the lemma, $s-\lambda(d)$ is strictly increasing in $d$ and changes sign exactly once, at $d^*$. Because $d_-$ is strictly decreasing in $K$, $h'>0$ for $K<K^*$ and $h'<0$ for $K>K^*$, where $d_-(K^*)=d^*$. The stated formula for $K^*$ inverts $d_-$. The final claim follows from the intermediate value theorem and strict monotonicity on each side.
\end{proof}

\begin{remark}
The market convention selects the preimage in $[K^*,\infty)$. At one month and $25\Delta$ this root lies far to the right of $K^*$; bracketing matters at long maturities and high volatilities.
\end{remark}

\section*{R6. Diffusive null of the hedge cost}

The comparison of hedged and unhedged carry returns follows Jurek (2014), who reads the resulting ratio as an upper bound on the crash share. The proposition below quantifies the part of that ratio that arises without any crash premium; the statement and proof are the author's.

Let the undiscounted put value with lognormal forward $G$ and volatility $\sigma$ be
\[
p(G)=K\,\Phi\bigl(-d_-(G)\bigr)-G\,\Phi\bigl(-d_+(G)\bigr),\qquad d_\pm(G)=\frac{\ln(G/K)\pm s^2/2}{s}.
\]
Using $K\varphi(d_-)=G\varphi(d_+)$, $p'(G)=-\Phi(-d_+(G))$.

Suppose that under $\Q$, $S_\tau$ is lognormal with mean $F$ and log-variance $s^2$, and that under $\P$ it is lognormal with the same log-variance and mean $Fe^{\lambda\tau}$ (no jumps, no volatility premium). A unit long forward costs nothing and has $\P$-expected payoff $\mu_U=F(e^{\lambda\tau}-1)$. Adding a put bought at its $\Q$-value, with the premium financed to $\tau$, gives $\mu_H=\mu_U+p(Fe^{\lambda\tau})-p(F)$.

\begin{proposition}
If $\lambda\ne0$, the hedge-cost ratio $\theta_0=1-\mu_H/\mu_U$ satisfies $\theta_0=\Phi(-d_+(\xi))$ for some $\xi$ between $F$ and $Fe^{\lambda\tau}$, and
\[
\bigl|\theta_0-\Phi(-d_1)\bigr|\le\varphi(0)\,\frac{|\lambda|\sqrt{\tau}}{\sigma},\qquad d_1=d_+(F).
\]
\end{proposition}

\begin{proof}
By the mean value theorem, $p(Fe^{\lambda\tau})-p(F)=-\Phi(-d_+(\xi))\,F(e^{\lambda\tau}-1)$ for some $\xi$ between the two forwards, so $\theta_0=-[p(Fe^{\lambda\tau})-p(F)]/\mu_U=\Phi(-d_+(\xi))$. Since $|d_+(\xi)-d_+(F)|=|\ln(\xi/F)|/s\le|\lambda|\tau/s$ and $\Phi$ is Lipschitz with constant $\varphi(0)$, the bound follows.
\end{proof}

\begin{remark}
With no crash premium at all, hedging with a put of forward delta $-\Phi(-d_1)$ surrenders approximately that fraction of the carry premium: about $0.25$ for a $25\Delta$ put and $0.10$ for a $10\Delta$ put. This is the null against which the realised ratio $\theta_{UB}$ and the skew component of E3 are read.
\end{remark}

\section*{R7. Entropy decomposition of currency premia}

The relation between the exchange rate and the two pricing kernels is that of Backus, Foresi and Telmer (2001). Conditional entropy, its cumulant expansion and the entropy bound are as in Backus, Chernov and Zin (2014), who relate the bound to Bansal and Lehmann (1997) and Alvarez and Jermann (2005). The combination below, including the covered-interest-parity term, is derived here.

Fix $t$ and write $\E_t$ for $\E[\cdot\mid\mathcal F_t]$. For a positive random variable $M$ with $\E_t|\log M|<\infty$, the conditional entropy is $L_t(M)=\log\E_tM-\E_t\log M\ge0$. Let $M,M^*>0$ be $\mathcal F_{t+1}$-measurable discount factors in $L^2$ that price domestic and foreign payoffs in their own currencies, with $\E_tM=e^{-r_t}$, $\E_tM^*=e^{-r^*_t}$, $\E_t|\log M|,\E_t|\log M^*|<\infty$ and $MS_{t+1}/S_t\in L^2$, where $S$ is the domestic price of foreign currency. Write $m=\log M$, $\Delta s=\log(S_{t+1}/S_t)$, and $x_t=(f_t-s_t)-(r_t-r^*_t)$ for the covered-interest-parity basis (on which see Du, Tepper and Verdelhan, 2018).

\begin{lemma}
Suppose every foreign payoff $X^*\in L^2(\mathcal F_{t+1})$ is traded and priced consistently in both currencies, $S_t\E_t[M^*X^*]=\E_t[MS_{t+1}X^*]$. Then $M^*=MS_{t+1}/S_t$ almost surely.
\end{lemma}

\begin{proof}
The difference $D=M^*-MS_{t+1}/S_t$ lies in $L^2$ and satisfies $\E_t[DX^*]=0$ for all $X^*\in L^2$. Taking $X^*=D$ gives $\E_tD^2=0$.
\end{proof}

\begin{proposition}
Under the lemma's conclusion, (i) $\E_t[\Delta s+r^*_t-r_t]=L_t(M)-L_t(M^*)$ and $\E_t[s_{t+1}-f_t]=L_t(M)-L_t(M^*)-x_t$. (ii) If $k(u)=\log\E_te^{um}$ is finite near $0$ and its Maclaurin series has radius of convergence greater than $1$, then $L_t(M)=\sum_{j\ge2}\kappa_{j,t}/j!$, where $\kappa_{j,t}$ are the conditional cumulants of $m$. (iii) For any gross return $R>0$ with $\E_t[MR]=1$ and $\E_t|\log R|<\infty$, $L_t(M)\ge\E_t\log R-r_t$.
\end{proposition}

\begin{proof}
(i) By definition, $\E_tm=-r_t-L_t(M)$ and $\E_tm^*=-r^*_t-L_t(M^*)$. The lemma gives $\Delta s=m^*-m$, so $\E_t\Delta s=r_t-r^*_t+L_t(M)-L_t(M^*)$. The second identity follows from $s_{t+1}-f_t=\Delta s-(f_t-s_t)=\Delta s-(r_t-r^*_t)-x_t$. (ii) $L_t(M)=k(1)-k'(0)$ with $k'(0)=\kappa_{1,t}$, and $k(1)=\sum_{j\ge1}\kappa_{j,t}/j!$ under the stated convergence. (iii) By Jensen's inequality, $0=\log\E_t[MR]\ge\E_tm+\E_t\log R$, so $L_t(M)=-r_t-\E_tm\ge\E_t\log R-r_t$.
\end{proof}

\begin{remark}
The identity in (i) needs only that some pair of discount factors prices bonds and exchange rates consistently; completeness makes the pair unique. Brandt, Cochrane and Santa-Clara (2006) show that the smoothness of exchange rates relative to the volatility of marginal utility implied by equity premia requires a very high degree of international risk sharing. Options on $S$ reveal the risk-neutral law of $\Delta s=m^*-m$, not the separate cumulants of $m$ and $m^*$, so they constrain but do not identify the split of $L_t(M)-L_t(M^*)$ across cumulant orders. Hedged and unhedged carry returns give two lower bounds on the same $L_t(M)$ through (iii).
\end{remark}

\section*{Results still to be written}

R3, R4, R5(c), R8 and the Garman--Kohlhagen PDE, as stated in \texttt{theory/README.md}.

\section*{References}
\begin{list}{}{\leftmargin=1.5em \itemindent=-1.5em \itemsep=2pt \parsep=0pt}\small
\item Alvarez, F. and Jermann, U. J. (2005). Using asset prices to measure the persistence of the marginal utility of wealth. \emph{Econometrica} 73(6), 1977--2016.
\item Backus, D., Chernov, M. and Zin, S. (2014). Sources of entropy in representative agent models. \emph{Journal of Finance} 69(1), 51--99.
\item Backus, D. K., Foresi, S. and Telmer, C. I. (2001). Affine term structure models and the forward premium anomaly. \emph{Journal of Finance} 56(1), 279--304.
\item Bansal, R. and Lehmann, B. N. (1997). Growth-optimal portfolio restrictions on asset pricing models. \emph{Macroeconomic Dynamics} 1(2), 333--354.
\item Brandt, M. W., Cochrane, J. H. and Santa-Clara, P. (2006). International risk sharing is better than you think, or exchange rates are too smooth. \emph{Journal of Monetary Economics} 53(4), 671--698.
\item Breeden, D. T. and Litzenberger, R. H. (1978). Prices of state-contingent claims implicit in option prices. \emph{Journal of Business} 51(4), 621--651.
\item Carr, P., Fisher, T. and Ruf, J. (2014). On the hedging of options on exploding exchange rates. \emph{Finance and Stochastics} 18(1), 115--144.
\item Du, W., Tepper, A. and Verdelhan, A. (2018). Deviations from covered interest rate parity. \emph{Journal of Finance} 73(3), 915--957.
\item Gatheral, J. and Jacquier, A. (2014). Arbitrage-free SVI volatility surfaces. \emph{Quantitative Finance} 14(1), 59--71.
\item Geman, H., El Karoui, N. and Rochet, J.-C. (1995). Changes of num\'eraire, changes of probability measure and option pricing. \emph{Journal of Applied Probability} 32(2), 443--458.
\item Jurek, J. W. (2014). Crash-neutral currency carry trades. \emph{Journal of Financial Economics} 113(3), 325--347.
\item Reiswich, D. and Wystup, U. (2012). FX volatility smile construction. \emph{Wilmott} 2012(60), 58--69. Version consulted: CPQF Working Paper No.~20, version 2, Frankfurt School of Finance \& Management, 20 March 2010.
\item Sampford, M. R. (1953). Some inequalities on Mill's ratio and related functions. \emph{Annals of Mathematical Statistics} 24(1), 130--132.
\item Shreve, S. E. (2004). \emph{Stochastic Calculus for Finance II: Continuous-Time Models}. Springer.
\end{list}

\end{document}
